\section{Modelos de Distribuição Discreta}

Muitas características de fenômenos aleatórios podem ser modelados
como variáveis aleatórias discretas, como determinar se o lançamento
de uma moeda deu cara, quantas pessoas possuem diploma numa população
de tamanho conhecido, quantas bolas devemos sortear de uma urna até
que obtenhamos a que contém a cor desejada, etc.

\subsection{Distribuição Bernoulli}

A distribuição Bernoulli modela experimentos que só podem ter dois resultados
possíveis: \textit{sucesso} ou \textit{fracasso}. Seu único parâmetro é a probabilidade $p$
de obter \textit{sucesso} no experimento. Dessa forma, podemos particionar o espaço
amostral em dois eventos disjuntos: $X = 0$ e $X = 1$. O primeiro é aquele cujos
resultados são \textit{fracassos}, obtidos com probabilidade $1 - p$, enquanto que no
segundo são aqueles cujos resultados são \textit{sucessos}, que ocorre com probablidade
$p$.

\begin{def:bernoulli}
Suponha que $\Omega$ seja um espaço amostral e $\mathcal{P} = \{F, S\}$ uma partição de
$\Omega$. A variável $X$ tem distribuição Bernoulli de parâmetro $p$, denotado por
$X \sim \mathrm{Bernoulli}(p)$, se, e somente se, $X$ é a função indicadora de um membro
de $\mathcal{P}$ e
\[
	P(X = 1) = p.
\]
\end{def:bernoulli}

Da definição formal de uma distribuição Bernoulli, segue que seu suporte é o conjunto
$\{0, 1\}$. Se $S \subset \Omega$ reúne os \textit{sucessos}, então podemos escrever
\[
	P(S) = p \quad \text{e} \quad P(F) = 1 - p
\]
e, caso $\Omega$ seja equiprovável, então
\[
	p = \frac{|S|}{|\Omega|}
\]

\begin{prop:bernoulli}
A f.m.p. da distribuição Bernoulli de parâmetro $p$ é dada por
\[
	p_X(x) = p^x(1-p)^{1-x}
\]
\begin{proof}
	Temos que
	\[
		P(X = x) = \begin{cases}
			p,   & \text{ se } x = 1 \\
			1-p, & \text{ se } x = 0 \\
		\end{cases}
	\]
	que equivale à fórmula desejada.
\end{proof}
\end{prop:bernoulli}

A esperança de uma Bernoulli é dada como
\[
	E[X] = 1 \cdot p + 0 \cdot (1-p) = p.
\]
Conclui-se que a taxa de sucessos é a própria probabilidade de obter um sucesso no
experimento.

\subsection{Distribuição Binomial}

A distribuição binomial é uma versão mais generalizada da Bernoulli, modelando quantos
sucessos são obtidos num experimento que consiste de $n$ tentativas Bernoulli
independentes, isto é, uma tentativa não altera no resultado da outra. Exemplo inclui
quantas caras obtemos ao lançar uma moeda 10 vezes. Os parâmetros dessa distribuição são
o número de tentativas $n$ e o número de sucessos $p$. Seu suporte será igual aos inteiros
de 0 a $n$.

Cada uma das $n$ tentativas Bernoulli são v.a.'s $X_k \sim \mathrm{Bernoulli}(p)$,
com $k \in [n]$, independentes. Logo, a probabilidade de um evento
$(X_k = x_k)_{k \in [n]}$ com $a$ sucessos, tendo fixado o valor de cada $x_k$,
deve ser igual a
\[
	P((X_k = x_k)_{k \in [n]}) = \prod_{k = 1}^n P(X_k = x_k),
\]
mas como há $a$ resultados $x_k = 1$, e $n-a$ resultados $x_k = 0$, então
\[
	P((X_k = x_k)_{k \in [n]}) = \prod_{k = 1}^n P(X_k = x_k) = p^{a} (1-p)^{n - a}
\]
Como cada uma das sequências de resultados com $a$ sucessos são mutalmente exclusivas,
então
\[
	P(X = a) = m \cdot p^{a} (1-p)^{n - a},
\]
onde $m$ é o número de sequências com $a$ sucessos dentre $n$, que equivale a
escolher $a$ tentativas para serem sucessos dentre $n$ possíveis, logo
\[
	P(X = a) = \binom{n}{a} \cdot p^{a} (1-p)^{n - a}
\]

\begin{def:binomial}
$
	X \sim \mathrm{Binomial}(n, p) \longleftrightarrow p_X(a) = \binom{n}{a} p^a (1-p)^{n-a}
$
\end{def:binomial}

Se a taxa de sucesso de cada ensaio Bernoulli é $p$, então podemos intuir que o número
esperado de sucessos em $n$ ensaios deverá ser a fração $p$ do total dos ensaios, ou seja
que o valor esperado seja $np$.

\begin{prop:esperanca binomial}
Se $X \sim \mathrm{Binomial}(n, p)$, então \(E[x] = np\).
\begin{proof}
	Com efeito,
	\[
		E[X] = \sum_{k=0}^n k \cdot \binom{n}{k} p^k (1-p)^{n-k}.
	\]
	Utilizando o fato de que
	\[
		\binom{n}{k} = \frac{n}{k} \binom{n-1}{k-1},
	\]
	reduzimos então a
	\[
		E[X] = n\sum_{k=0}^n \binom{n-1}{k-1} p^k (1-p)^{n-k},
	\]
	o que nos sugere realizar uma troca de variável em $k$ para que consigamos usar o
	Teorema Binomial. Para tanto, realizemos que
	\[
		n\sum_{k=0}^n \binom{n-1}{k-1} p^k (1-p)^{n-k} = np \sum_{k=1}^n\binom{n-1}{k-1}p^{k-1}(1-p)^{n-k}
	\]
	já que somar o termo $k=0$ equivale a somar 0. Fazendo então a troca $j = k - 1$, obtemos
	\[
		E[X] = np \sum_{j=0}^n\binom{n-1}{j}p^{j}(1-p)^{(n-1)-j},
	\]
	e pelo Teorema Binomial,
	\[
		E[X] = np \cdot (p + 1 - p)^{n-1} = np
	\]
\end{proof}
\end{prop:esperanca binomial}

\subsection{Distribuição Geométrica}

Uma variável geométrica modela experimentos que consistem na repetição de ensaios Bernoulli
independentes até que se obtenha o primeiro sucesso. Um exemplo é repetir o lançamento
de uma moeda até obter coroa. Para obter a função de massa de probabilidade dessa variável
para uma realização $n$, basta encontrar a probabilidade de obter $n-1$ fracassos
consecutivos seguidos de um sucesso. Desse modo, se $(X_k = x_k)_{k \in [n]}$ é
uma sucessão de $n$ ensaios Bernoulli independentes com parâmetro $p$, então
\[
	P((X_k = x_k)_{k \in \mathbb{N}}) = \prod_{k=1}^n P(X_k = x_k),
\]
mas como os $n-1$ valores de $x_k$ são fracassos (iguais a 0) obtidos com probabilidade
$1-p$, e $x_n = 1$ é obtido com probabilidade $p$, então
\[
	P((X_k = x_k)_{k \in \mathbb{N}}) = (1-p)^{n-1}p
\]

\begin{def:geometrica}
\[
	X \sim Geo(p) \longleftrightarrow p_X(n) = (1-p)^{n-1}p
\]
\end{def:geometrica}

O número esperada de ensaios realizados num experimento até obter o primeiro sucesso
é o valor esperado da variável geométrica.

\begin{prop:esperanca geometrica}
Se $X \sim \mathrm{Geo}(p)$, então
\[
	E[X] = \frac{1}{p}
\]
\begin{proof}
	Com efeito,
	\[
		E[X] = \sum_{k=0}^\infty k (1-p)^{k-1} p = p \sum_{k=0}^\infty k(1-p)^{k-1}.
	\]
	O somatório à esquerda é a derivada de uma série geométrica de razão $q = 1-p$,
	tal que $q < 1$. Observemos que
	\begin{align*}
		\frac{d}{dq} \left(\sum_{k=0}^\infty q^k\right) & =
		\frac{d}{dq} \left(\frac{1}{1 - q}\right)                             \\
		\sum_{k=0}^\infty k(1-p)^{k-1}                  & = \frac{1}{(1-q)^2}
	\end{align*}
	donde conclui-se
	\[
		E[X] = \frac{p}{(1 - 1 + p)^2} = \frac{1}{p}
	\]
\end{proof}
\end{prop:esperanca geometrica}

\subsection{Distribuição Binomial Negativa}

A distribuição binomial negativa é similar à geométrica: invés de repetir os ensaios
Bernoulli independentes até o obter o primeiro sucesso, os ensaios são repetidos até
que se acumulem $r$ sucessos, sendo esse total um novo parâmetro da distribuição.
A dedução da f.m.p. para esse modelo de probabilidade é análogo ao da Binomial e da
Geométrica; fixando uma sequência $((X_k = x_k)_{k \in [n]}$ de ensaios Bernoulli
independentes com parâmetro $p$, temos que a probabilidade de uma sequência será
\[
	\prod_{k=1}^n P(X_k = x_k),
\]
de sorte que, sendo $r$ ensaios com $x_k = 1$ e $n-r$ ensaios com $x_k = 0$, então
\[
	\prod_{k=1}^n P(X_k = x_k) = (1-p)^{n-r} p^r.
\]
Visto que cada sequência de ensaios é mutualmente exclusiva, então a probabilidade
de obter acumular $n$ ensaios até obter $r$ tentativas equivale ao total de escolhas
de $r$ ensaios para fixar $x_k = 1$ dentre $n$ possíveis, logo
\[
	\binom{n}{r} (1-p)^{n-r} p^r
\]

\begin{def:binomial negativa}
\[
	X \sim \mathrm{BinNeg}(r, p) \longleftrightarrow p_X(n) = \binom{n}{r} (1-p)^{n-r} p^r
\]
\end{def:binomial negativa}

Uma binomial negativa pode ser vista como a sucessão de várias geométricas. Desse modo,
o valor esperado de uma binomial negativa deve ser igual ao valor esperado da soma das
geométricas que a compõe.

\begin{prop:esperanca binomial negativa}
Se $X \sim \mathrm{BinNeg}(r, p)$, então
\[
	E[X] = \frac{r}{p}
\]
\begin{proof}
	Suponha que $(X_i)_{i \in [r]}$ denote a variável que determina quantos
	ensaios Bernoulli foram realizados ate se obter um sucesso desde o último. Com efeito,
	a variável $X_i \sim \mathrm{Geo}(p)$ para todo $i \in [r]$. Além disso, temos que
	\[
		X = \sum_{i = 1}^r X_i
	\]
	Pela linearidade da esperança, conclui-se que
	\[
		E[X] = \sum_{i = 1}^r E[X_i] = \sum_{i=1}^r \frac{1}{p} = \frac{r}{p}
	\]
\end{proof}
\end{prop:esperanca binomial negativa}

\subsection{Distribuição Hipergeométrica}

A distribuição hipergeométrica modela amostragens de tamanho $n$ de um espaço amostral
equiprovável de tamanho $N$, em que $m$ dos resultados possíveis são considerados
sucessos. Um exemplo clássico são bolas numa urna, das quais $m$ são azuis, e as demais
vermelhas. A urna é nosso espaço amostral, as $m$ bolas azuis são os resultados que
representam sucessos possíveis, caso realizassemos um sorteio de $n$ bolas da urna,
a variável que nos dá o número de bolas azuis sorteadas é hipergeométrica.

Para obter a função de distribuição para a situação sendo modelada, considere uma população
de tamanho $N$, com $m$ indivíduos pertencentes ao evento de sucessos, e uma amostra de
tamanho $n$ dessa população. Para obter a probabilidade de obter $k$ indivíduos na amostra
que sejam sucessos, supondo que os indivíduos na população são equiprováveis, podemos
utilizar a proposição \ref{prop:equiprovavel}, considerando $\Omega$ como o universo das
amostras, e $E$ as amostras que possuem $k$ indivíduos considerados sucessos. O número de
amostras possíveis equivale a selecionar $n$ indivíduos de $N$ possíveis
\[
	|\Omega| = \binom{N}{n}.
\]

Por sua vez, o número de amostras com $k$ sucessos equivale ao selecionar $k$ indivíduos
dos $m$ que são sucessos e selecionar outros $n-k$ com que não pertencem a esse grupo dos
$N-m$ restantes. Por conseguinte
\[
	|E| = \binom{m}{k} \cdot \binom{N-m}{n-k}
\]
e a probabilidade de uma amostra conter $k$ sucessos será
\[
	\frac{\binom{m}{k} \cdot \binom{N-m}{n-k}}{\binom{N}{n}}
\]

\begin{def:hipergeometrica}
\[
	X \sim \mathrm{HG}(N, m, n) \longleftrightarrow p_X(k) =
	\frac{\binom{m}{k} \cdot \binom{N-m}{n-k}}{\binom{N}{n}}
\]
\end{def:hipergeometrica}

A média de uma hipergeométrica pode ser facilmente obtida pelo seguinte
raciocínio: suponha que já tenhamos selecionado a amostra de tamanho $n$ e temos
o conhecimento \textit{a priori} de que a população $N$ continha $m$ sucessos,
implicando que a probabilidade que cada indivíduo da amostra selecionada tenha
probabilidade $m / N$ de ser um sucesso. Considerando que cada indivíduo na
amostra é um ensaio Bernoulli de parâmetro $m / N$, então a média de sucessos
deve ser igual a $\frac{nm}{N}$.

\begin{prop:esperanca hipergeometrica}
Se $X \sim \mathrm{HG}(N, m, n)$, então
\[
	E[X] = n \frac{m}{N}
\]
\begin{proof}
	Presuma que a amostra de tamanho $n$ obtida seja a sequência de eventos
	$(X_k = x_k)_{k \in [n]}$, em que $X_k: \Omega \to \{0, 1\}$ é a variável indicadora
	do conjunto dos sucessos $S$. Como $|S| = m$ e $|\Omega| = N$, então
	\[
		P(S) = \frac{m}{N}
	\]
	e $X_k$ equivale a uma v.a. Bernoulli de parâmetro $m / N$. Logo, o número de
	sucessos na amostra é uma variável binomial com parâmetros $n$ e $n / N$, e sua
	média será
	\[
		E[X] = n \frac{m}{N}
	\]
\end{proof}
\end{prop:esperanca hipergeometrica}

\subsection{Distribuição de Poisson}

A distrbuição de Poisson modela contagem de eventos que ocorrem dentro de um intervalo
contínuo -- como o tempo, uma superfície ou volume -- com probabilidade praticamente
desprezível. Essa distribuição pode ser compreendida como a uma aproximação da Binomial
quando $n \to \infty$.

Considere que $X \sim \mathrm{Binomial}(n, p)$, logo
\[
	p_X(a) = \binom{n}{a} (1-p)^{n-a} p^a
\]
Se tomarmos o limite da f.m.p. de $X$ quando $n \to \infty$ e denotamos $np$ por $\lambda$,
então
\begin{align*}
	\lim_{n \to \infty} p_X(a) & = \lim_{n \to \infty} \binom{n}{a} \left(\frac{\lambda}{n}\right)^a \left(1 - \frac{\lambda}{n}\right)^{n-a}                  \\
	                           & = \lim_{n \to \infty} \frac{n(n-1)\ldots(n-a+1)}{a!}\left(\frac{\lambda}{n}\right)^a \left(1 - \frac{\lambda}{n}\right)^{n-a} \\
	                           & = \lim_{n \to \infty} \prod_{k=1}^{a-1} \left(1 - \frac{k}{n}\right) \cdot
	\lim_{n \to \infty} \frac{\lambda^a}{a!} \frac{(1 - \lambda / n)^n}{(1 - \lambda / n)^a}
\end{align*}
Como
\[
	\lim_{n \to \infty} \prod_{k=1}^{a-1} \left(1 - \frac{k}{n}\right) = 1 \quad
	\text{e} \quad \lim_{n \to \infty} 1 - (\lambda / n)^a = 1
\]
então,
\[
	\lim_{n \to \infty} p_X(a) = \lim_{n \to \infty} \frac{\lambda^a}{a!} \left(1 - \frac{\lambda}{n}\right)^n \\
\]
Usando o limite exponencial fundamental
\[
	\lim_{n \to \infty} \left(1 - \frac{\lambda}{n}\right)^n = e^{-\lambda}
\]
obtemos, por fim,
\[
	\lim_{n \to \infty} p_X(a) = e^{-\lambda} \frac{\lambda^a}{a!}
\]

\begin{def:poisson}
\[
	X \sim \mathrm{Poisson}(\lambda) \longleftrightarrow p_X(a) = e^{-\lambda} \frac{\lambda^a}{a!}
\]
\end{def:poisson}

\begin{prop:esperanca poisson}
Se $X \sim \mathrm{Poisson}(\lambda)$, então
\[
	E[X] = \lambda
\]
\begin{proof}
	Com efeito,
	\begin{align*}
		E[X] & = \sum_{k = 0}^\infty k e^{-\lambda} \frac{\lambda^k}{k!}             \\
		     & = \lambda e^{-\lambda} \sum_{k=1}^\infty \frac{\lambda^{k-1}}{(k-1)!}
	\end{align*}
	O somatório equivale ao Polinômio de Taylor de $e^\lambda$, portanto
	\[
		E[X] = \lambda e^{-\lambda} e^{\lambda} = \lambda
	\]
\end{proof}
\end{prop:esperanca poisson}
